%% 4. なぜ効くか（実験 3 と探索）
\setband{\T{\S4 なぜ効くか}{\S4 Why it works}}

\begin{frame}{\T{得が出るのは，教師が生徒より当たっているときである}{Gains appear where the teacher is more accurate than the student}}
\rl{\T{探索（評価の後）}{Exploration (post-evaluation)}}
\hosoku{appadv}{\T{K = 25 の図と数値}{chart and numbers for K = 25}}
\begin{columns}[T]
\begin{column}{0.54\textwidth}
\begin{tikzpicture}
\begin{axis}[width=\linewidth, height=6.2cm, xmin=-2, xmax=33, ymin=-3.5, ymax=9.5,
  xlabel={\T{教師の先回り（候補での UAR，教師 − 生徒．ポイント）}{teacher lead (pool UAR, teacher − student, points)}},
  ylabel={\T{提案の得（テスト UAR，提案 − K 本だけ）}{gain (test UAR, proposal − K-only)}},
  legend style={at={(0.98,0.03)}, anchor=south east, font=\tiny, cells={anchor=west}}, legend columns=2, grid=major,
  every node near coord/.append style={font=\tiny, anchor=west, xshift=2pt, color=mMuted}]
\addplot[only marks, kA, mark=*, mark size=2pt, point meta=explicit symbolic, nodes near coords] table[x=adv_pool, y=gain_prop, meta=T]{teacher_adv_A_K5.csv};
\addplot[only marks, kV, mark=square*, mark size=2pt, point meta=explicit symbolic, nodes near coords] table[x=adv_pool, y=gain_prop, meta=T]{teacher_adv_V_K5.csv};
\addplot[only marks, kT, mark=triangle*, mark size=2.4pt, point meta=explicit symbolic, nodes near coords] table[x=adv_pool, y=gain_prop, meta=T]{teacher_adv_T_K5.csv};
\addplot[only marks, kAVT, mark=diamond*, mark size=2.4pt] table[x=adv_pool, y=gain_prop]{teacher_adv_AVT_K5.csv};
\legend{\T{音声}{audio}, \T{顔}{face}, \T{言葉}{text}, AVT}
\end{axis}
\end{tikzpicture}
\end{column}
\begin{column}{0.46\textwidth}
\footnotesize
\T{%
\begin{itemize}
\item 横軸：教師（3 モダリティの積）が候補でどれだけ生徒（K 本だけ）より当たるか．縦軸：提案のテストでの得．K = 5，16 点（4 コーパス × 4 生徒）
\item 順位の相関は \N{rhofive}（K = 25 で \N{rhotwofive}）．先回りが大きい弱いモダリティ（音声，IEMOCAP の言葉）に得が出る
\item \textbf{AVT には学ぶものが無い}（図の菱形 4 点，コーパス名なし）：先回りは \N{advMOSEIAVTfive}〜\N{advAFEWAVTfive} ポイントで，教師と AVT の生徒の判定は候補の \N{agreeMOSEIAVTfive}〜\N{agreeAFEWAVTfive}\% で一致する
\end{itemize}
}{%
\begin{itemize}
\item x: how much the teacher (product of 3 modalities) beats the student (K-only) on the pool. y: the proposal's test gain. K = 5, 16 points (4 corpora × 4 students)
\item Rank correlation \N{rhofive} (\N{rhotwofive} at K = 25). Gains appear for weak modalities with a large lead (audio, IEMOCAP text)
\item \textbf{Nothing to learn for AVT} (the 4 unlabeled diamonds): the lead is \N{advMOSEIAVTfive} to \N{advAFEWAVTfive} points, and the teacher agrees with the AVT student on \N{agreeMOSEIAVTfive}--\N{agreeAFEWAVTfive}\% of the pool
\end{itemize}
}
\end{column}
\end{columns}
{\tiny\color{mMuted}\T{評価の回を見た後の探索．16 点は独立でなく，p 値は付けない．MOSEI は途中の 39 単位．出どころ：figures/teacher\_adv\_*.csv}{Post-evaluation exploration; the 16 points are not independent and no p-value is attached. MOSEI: 39 interim units. Source: figures/teacher\_adv\_*.csv}}
\end{frame}

\begin{frame}{\T{損を「位置」と「向き」に分ける：出発点の損の大半は重みの向きのずれである}{Splitting the loss into ``position'' and ``direction'': most of the start's loss is a weight-direction error}}
\rl{\T{実験 3}{Exp.~3}}
\hosoku{appdecomp}{\T{全生徒・全 K の表}{tables for all students and K}}
\begin{columns}[T]
\begin{column}{0.5\textwidth}
\centering
\begin{tikzpicture}[font=\scriptsize, x=1cm, y=1cm]
  \pgfmathsetseed{7}
  % 2 つの雲
  \foreach \i in {1,...,14} {\pgfmathsetmacro{\px}{0.9+0.75*rand}\pgfmathsetmacro{\py}{1.7+0.5*rand}\fill[kA!70] (\px,\py) circle (1.3pt);}
  \foreach \i in {1,...,14} {\pgfmathsetmacro{\px}{2.6+0.75*rand}\pgfmathsetmacro{\py}{0.5+0.5*rand}\fill[kV!70] (\px,\py) circle (1.3pt);}
  \node[font=\tiny, kA] at (0.9,2.45) {\T{感情 1}{emotion 1}}; \node[font=\tiny, kV] at (2.6,-0.25) {\T{感情 2}{emotion 2}};
  % 正しい境界
  \draw[cUpper, line width=1pt] (0.2,-0.2) -- (3.4,2.6);
  % 位置のずれ（平行移動）
  \draw[cCrit, line width=1pt, densely dashed] (0.9,-0.2) -- (4.1,2.6);
  % 向きのずれ（回転）
  \draw[cDisc, line width=1pt, dotted] (-0.4,0.8) -- (4.0,1.6);
  \node[font=\tiny, anchor=north, align=left] at (1.9,-0.55) {\textcolor{cUpper}{\rule[0.5ex]{2em}{1pt}} \T{正しい境界}{right boundary}\quad \textcolor{cCrit}{- - -} \T{位置のずれ（切片：平行に動く）}{position error (intercept: parallel shift)}\\ \textcolor{cDisc}{$\cdots$} \T{向きのずれ（重み：回る）}{direction error (weights: rotation)}};
\end{tikzpicture}
\notreal
\vspace{0.2em}
\scriptsize
\T{測り方：識別器の切片だけをテストの感情の釣り合いで当て直す．当て直しで減る分＝位置の損，残る分＝向きの損（上限との差）．}{How measured: refit only the intercept on the test emotion balance. The reduction = position loss; what remains = direction loss (gap to the upper bound).}
\end{column}
\begin{column}{0.5\textwidth}
\begin{tikzpicture}
\begin{axis}[width=\linewidth, height=4.4cm, ybar stacked, bar width=14pt, ymin=0, ymax=18, enlarge x limits=0.2,
  xtick={1,2,3,4}, xticklabels={AFEW, IEMOCAP, MELD, MOSEI}, ylabel={\T{出発点の損（AVT，ポイント）}{loss of the start (AVT, points)}},
  legend style={at={(0.98,0.97)}, anchor=north east, font=\tiny, cells={anchor=west}}, ymajorgrids,
  nodes near coords={\pgfmathprintnumber[fixed, precision=1]{\pgfplotspointmeta}}, every node near coord/.append style={font=\tiny, color=white}]
\addplot[fill=cCrit, draw=cCrit] table[x=ci, y=position]{decomp_start_avt.csv};
\addplot[fill=cDisc, draw=cDisc] table[x=ci, y=direction]{decomp_start_avt.csv};
\legend{\T{位置（切片）}{position (intercept)}, \T{向き（重み）}{direction (weights)}}
\end{axis}
\end{tikzpicture}
\footnotesize
\T{%
\begin{itemize}
\item AVT の出発点：向きの損 \N{dirMOSEIAVTstart}〜\N{dirAFEWAVTstart}，位置の損 \N{posMOSEIAVTstart}〜\N{posIEMOCAPAVTstart} ポイント
\item 向きの物差しの確かめに不合格の組（AFEW と IEMOCAP の AVT など 5 組）では，論文で「向き」と呼ばず「当て直した後に残る差」と書く
\end{itemize}
}{%
\begin{itemize}
\item AVT start: direction loss \N{dirMOSEIAVTstart} to \N{dirAFEWAVTstart}, position loss \N{posMOSEIAVTstart} to \N{posIEMOCAPAVTstart} points
\item Where the direction measure failed its validity check (5 pairs incl.\ AFEW and IEMOCAP AVT), the paper says ``gap remaining after the refit'' instead of ``direction''
\end{itemize}
}
\end{column}
\end{columns}
\end{frame}

\begin{frame}{\T{K 本も向きを直していた．提案は，弱いモダリティの向きの損をさらに減らす}{K labels also fix the direction; the proposal further shrinks the direction loss of weak modalities}}
\rl{\T{実験 3}{Exp.~3}}
\begin{columns}[T]
\begin{column}{0.56\textwidth}
\footnotesize
\T{向きの損（当て直した後に残る差，ポイント，K = 5）}{Direction loss (gap remaining after the refit, points, K = 5)}
\begin{tabular}{@{}llrrr@{}}
\toprule
\T{コーパス}{Corpus} & \T{生徒}{Student} & \T{出発点}{Start} & \T{K 本だけ}{K-only} & \T{提案}{Proposal} \\
\midrule
AFEW & \T{音声}{audio} & \N{dirAFEWAstart} & \N{dirAFEWAnonefive} & \textbf{\N{dirAFEWApropfive}} \\
IEMOCAP & \T{音声}{audio} & \N{dirIEMOCAPAstart} & \N{dirIEMOCAPAnonefive} & \textbf{\N{dirIEMOCAPApropfive}} \\
MELD & \T{音声}{audio} & \N{dirMELDAstart} & \N{dirMELDAnonefive} & \textbf{\N{dirMELDApropfive}} \\
MOSEI & \T{音声}{audio} & \N{dirMOSEIAstart} & \N{dirMOSEIAnonefive} & \textbf{\N{dirMOSEIApropfive}} \\
IEMOCAP & \T{言葉}{text} & \N{dirIEMOCAPTstart} & \N{dirIEMOCAPTnonefive} & \textbf{\N{dirIEMOCAPTpropfive}} \\
MELD & \T{言葉}{text} & \N{dirMELDTstart} & \N{dirMELDTnonefive} & \N{dirMELDTpropfive} \\
MOSEI & \T{言葉}{text} & \N{dirMOSEITstart} & \N{dirMOSEITnonefive} & \N{dirMOSEITpropfive} \\
MELD & AVT & \N{dirMELDAVTstart} & \N{dirMELDAVTnonefive} & \N{dirMELDAVTpropfive} \\
MOSEI & AVT & \N{dirMOSEIAVTstart} & \N{dirMOSEIAVTnonefive} & \N{dirMOSEIAVTpropfive} \\
\bottomrule
\end{tabular}
\end{column}
\begin{column}{0.44\textwidth}
\footnotesize
\T{%
\begin{itemize}
\item 事前に立てた説明「K 本は主に切片を直す」は 4 コーパスとも支持されない．K 本は向きも直す（例：IEMOCAP の言葉 \N{dirIEMOCAPTstart} → \N{dirIEMOCAPTnonefive}）
\item 提案は，得の出る生徒で向きの損をさらに減らす（IEMOCAP の言葉 \N{dirIEMOCAPTnonefive} → \N{dirIEMOCAPTpropfive}，音声は 4 コーパスで減る）
\item 得の出ない生徒では減らない，または増える（MELD・MOSEI の言葉，MOSEI の AVT）
\item 「重みは予測で，切片は K 本で」と役割を分けた版は，同時に学ぶ版と差が無い（−0.4〜+1.7）
\end{itemize}
}{%
\begin{itemize}
\item The pre-registered ``K labels mainly fix the intercept'' is not supported in any corpus; K labels also fix the direction (e.g.\ IEMOCAP text \N{dirIEMOCAPTstart} → \N{dirIEMOCAPTnonefive})
\item The proposal further reduces the direction loss where gains appear (IEMOCAP text \N{dirIEMOCAPTnonefive} → \N{dirIEMOCAPTpropfive}; audio in all 4 corpora)
\item Where there is no gain, it does not shrink or grows (MELD and MOSEI text, MOSEI AVT)
\item A version that learns weights from predictions and the intercept from K labels does not differ from learning both together (−0.4 to +1.7)
\end{itemize}
}
\end{column}
\end{columns}
{\tiny\color{mMuted}\T{点の値．区間と判定の規則は補足．出どころ：figures/decomp.csv（版 four）}{Point values; intervals and decision rules in the appendix. Source: figures/decomp.csv (version four)}}
\end{frame}

\begin{frame}{\T{効かない所の理由：顔は学ぶものが無いか切片が悪くなる．MELD・MOSEI は教師そのものが弱い}{Why it fails: the face has nothing to learn or gets a worse intercept; in MELD and MOSEI the teacher itself is weak}}
\rl{\T{実験 3・探索}{Exp.~3, exploration}}
\begin{columns}[T]
\begin{column}{0.5\textwidth}
\footnotesize
\textbf{\T{顔の生徒}{Face student}}
\T{%
\begin{itemize}
\item AFEW：教師の先回りは \N{advAFEWVfive} ポイントしかない（顔が最も強いモダリティ）→ 学ぶものが無い
\item IEMOCAP：損の大半が位置（切片）で（\N{posIEMOCAPVnonefive} 対 向き \N{dirIEMOCAPVnonefive}），提案は位置の損を \N{posIEMOCAPVnonefive} → \N{posIEMOCAPVpropfive} と\textbf{悪くする}
\item MELD：先回り \N{advMELDVfive} に対して得は \N{gainrMELDVfive}．理由は未解析
\item 直す余地（上限 − K 本だけ）は顔で \N{roomMELDVfive}〜\N{roomMOSEIVfive}，音声で \N{roomAFEWAfive}〜\N{roomMOSEIAfive} ポイント
\end{itemize}
}{%
\begin{itemize}
\item AFEW: the teacher's lead is only \N{advAFEWVfive} points (the face is the strongest modality) → nothing to learn
\item IEMOCAP: most of the loss is position (intercept; \N{posIEMOCAPVnonefive} vs direction \N{dirIEMOCAPVnonefive}), and the proposal \textbf{worsens} it (\N{posIEMOCAPVnonefive} → \N{posIEMOCAPVpropfive})
\item MELD: lead \N{advMELDVfive} but gain \N{gainrMELDVfive}; not yet analyzed
\item Room (upper bound − K-only): face \N{roomMELDVfive}--\N{roomMOSEIVfive}, audio \N{roomAFEWAfive}--\N{roomMOSEIAfive} points
\end{itemize}
}
\end{column}
\begin{column}{0.5\textwidth}
\footnotesize
\textbf{\T{教師の強さ（候補での UAR，K = 5）}{Teacher strength (pool UAR, K = 5)}}
\begin{tabular}{@{}lrr@{}}
\toprule
\T{コーパス}{Corpus} & \T{教師}{Teacher} & \T{AVT の生徒（K 本だけ）}{AVT student (K-only)} \\
\midrule
AFEW & \N{tuarAFEWAVTfive} & \N{suarAFEWAVTfive} \\
IEMOCAP & \N{tuarIEMOCAPAVTfive} & \N{suarIEMOCAPAVTfive} \\
MELD & \N{tuarMELDAVTfive} & \N{suarMELDAVTfive} \\
MOSEI & \N{tuarMOSEIAVTfive} & \N{suarMOSEIAVTfive} \\
\bottomrule
\end{tabular}
\T{%
\begin{itemize}
\item MELD と MOSEI では教師が候補の半分も当てない．それでも音声には効く（音声の生徒がもっと弱いから）
\item 教師の強さはコーパスで大きく違う．\textbf{どのモダリティをどれだけ信じるか}がコーパスごとに違うことが，次の問題になる（\S5，\S6）
\end{itemize}
}{%
\begin{itemize}
\item In MELD and MOSEI the teacher gets less than half the pool right, yet it still helps audio (the audio student is weaker still)
\item Teacher strength differs widely across corpora. \textbf{Which modality to trust, and how much}, differs per corpus --- the next problem (\S5, \S6)
\end{itemize}
}
\end{column}
\end{columns}
\end{frame}
